\section{The stationary adjoint at nondegenerate equality}
\label{j:section}

Assume equality in Theorem~\ref{m:main} and the strict area branch. Throughout
this section, \(g,k\) denote the \emph{polarized} metric and tensor of
Equation~\eqref{a:constraints}. Write
\[
 R_0=\sqrt{\frac{A}{4\pi}},\qquad M=m,\qquad
 \kappa=\frac{b}{R_0^2}>0,
\]
where the subextremal parameters are those of
Lemma~\ref{a:model-parameters}, evaluated at \(R=R_0\).
Thus \(g=h+X^2\dd\phi^2\), and \(H+\tr_Sk=0\).
We write \(D_i\) for the three rows of \(\cD_X\), and set
\[
 (h_1,h_2,h_3)=(1,1,2),\qquad e=X^{-1}\eta.
\]
The numbers \(h_i\) here are weights, distinct from the meridional metric \(h\).

\begin{proposition}[Stationary adjoint]
\label{j:adjoint}
There are smooth invariant fields \(N,Y\) on the polarized exterior,
with \(Y\) meridional, such that
\begin{equation}
\label{j:causal}
 N\geq |Y|_g,\qquad N>0\ \text{in }\operatorname{int}\Omega,\qquad
 \operatorname{sym}\nabla Y=-Nk.
\end{equation}
For every \(q'\in(1/2,1)\), their asymptotics are
\begin{equation}
\label{j:asymptotic}
 N=1+O_2(r^{-q'}),\qquad Y=O_2(r^{-2}).
\end{equation}
At the inner boundary, with \(\nu\) directed into the exterior,
\begin{equation}
\label{j:boundary}
 Y=N\nu,\qquad \partial_\nu N+k(Y,\nu)=\kappa.
\end{equation}
Furthermore, either \(N|_S>0\) everywhere or \(N|_S\equiv0\).

On regular orbits in the interior, define
\begin{equation}
\label{j:stationary-metric}
 H_3=-N^2\dd t_0^2+
       h_{ij}(\dd x^i+Y^i\dd t_0)(\dd x^j+Y^j\dd t_0),
 \qquad \gamma=X^2H_3.
\end{equation}
Then \((\gamma,\Phi)\), with \(\partial_{t_0}\Phi=0\), solves the full
Einstein--wave-map system of Equation~\eqref{a:wave}. For the future
\(H_3\)-unit normal \(n\) to \(t_0=0\), the induced horizontal tensor is
\(k|_h\), and the map velocities are
\[
 \dd\log X(n)=k(e,e),\qquad \cD_X(n)=s.
\]
\end{proposition}

The proof uses the numerical inequality on its enlarged sourced-constraint
class, as established in Corollary~\ref{d:numerical}. The equality-variation
method has a neutral antecedent in \cite{NeutralExterior2026}; all the
constraint and matter variations needed here are given below.

\subsection{The minimum-area envelope}

For a polarized metric \(\widetilde g\) near \(g\), let \(a(\widetilde g)\)
be the minimum enclosing area. The argument uses its one-sided derivative;
it does not require uniqueness of a minimizing enclosure.

\begin{lemma}[Compact minimizing family and envelope derivative]
\label{j:area}
The infimum defining \(a(\widetilde g)\) is attained by an invariant
filled inner set. The family \(\cT_0\) of invariant minimizers at \(g\),
viewed through their sets modulo null sets and their area and tangent-plane
measures, is compact. At the equality data,
\[
 a(g)=A.
\]
For any of the smooth metric paths \(\widetilde g_d\) used below,
\begin{equation}
\label{j:envelope}
 \left.\frac{\dd}{\dd d}\right|_{0+}a(\widetilde g_d)
   =\min_{T\in\cT_0} A'_T,\qquad
 A'_T=\frac12\int_{\partial T}\tr_{\partial T}\dot g\,\dd A_g.
\end{equation}
\end{lemma}

\begin{proof}
Extend the geometry smoothly and equivariantly across \(S\) into an
auxiliary cap, and require each bounded inner set to contain that cap up
to \(S\). Count its full frontier, including coincidence with the obstacle.
First impose an outer sphere obstacle. Perimeter compactness gives a
minimizer. It has a \(C^{1,1}\) frontier, smooth and minimal off the inner
obstacle, by smooth-obstacle perimeter regularity in dimension three
\cite[Regularity Theorem~1.3(iii)]{HuiskenIlmanen2001}.

The outer sphere can be removed uniformly for metrics in the paths under
consideration. Large coordinate spheres are strictly mean convex, and the
competitor \(S\) gives a uniform area upper bound. If a free minimizing
frontier had a point at radius \(r_j\to\infty\), its portion in an
asymptotic ball of radius \(\varepsilon r_j\), for a fixed small
\(\varepsilon>0\), would have area at least \(c r_j^2\) by the interior
density estimate for perimeter minimizers. This contradicts the area upper
bound. Thus every minimizer stays in one compact set. Filling bounded
pockets decreases perimeter and preserves admissibility.

Among minimizers in this compact set choose one with largest filled
volume. The perimeter inequality for union and intersection shows that
both the union and the intersection of two minimizing inner sets are
again minimizers. Rotating a maximum-volume minimizer and taking its
union with the original one therefore leaves it unchanged modulo null
sets. It is invariant. Smooth outward approximation of its regular
graphs, or equivariant smoothing of an outward defining function, gives
smooth invariant enclosing cuts with convergent area. Invariant areas are
preserved by polarization. The original minimum-area hypothesis consequently
gives \(a(g)\geq A\), while \(S\) gives the reverse inequality.

For a sequence \(d_j\to0\), choose invariant minimizers \(T_j\) for
\(\widetilde g_{d_j}\). Compact containment and perimeter compactness give
a subsequence converging to a filled inner set \(T\). Uniform equivalence
of the metrics, lower semicontinuity, and comparison with a background
minimizer show that \(T\in\cT_0\) and that the perimeters converge.
Strict convergence of the vector measures representing the perimeter
normals then implies convergence of their area and tangent-plane
measures. In particular, the integral defining \(A'_T\) is continuous
under this convergence.

On the common compact set the expansion of area is uniform:
\[
 \Area_{\widetilde g_d}(\partial T)
 =\Area_g(\partial T)+dA'_T+O(d^2)\Area_g(\partial T).
\]
Comparing the perturbed minimum first with a minimizing member of
\(\cT_0\), and then with its own minimizer \(T_d\), proves the upper and
lower one-sided bounds in Equation~\eqref{j:envelope}. The same compactness
argument at \(d=0\) proves compactness of \(\cT_0\).
\end{proof}

We will also use the common tangent plane of minimizing frontiers wherever
they meet off \(S\). A transverse meeting would give a corner in the
minimizing union or intersection, contradicting its regularity. This
common plane is a measurable function on the union of the free frontiers.
Here is a compact-graph verification, including accumulation leaves.
For \(\delta_1>0\), consider the triples \((T,x,\Pi)\) with
\(T\in\cT_0\), \(x\in\partial T\),
\(\dist(x,S)\geq\delta_1\), and \(\Pi=T_x\partial T\).
Their set is compact. Indeed the density bound prevents a limiting support
point from disappearing, while strict perimeter convergence and free
minimizing-boundary regularity give graphical \(C^1\) convergence near
the smooth multiplicity-one limiting frontier. Thus the limiting plane
is its tangent plane. Projection onto \((x,\Pi)\) is a compact graph;
the common-plane property makes it single-valued. Its plane function is
therefore continuous on its compact projection. Exhausting by
\(\delta_1\downarrow0\) gives the required measurable plane field.

An invariant crossing arc gives an enclosure after rotation and smooth
simple crossing approximation. Its length in \(X^2h\) is therefore at
least \(a(\widetilde g)/(2\pi)\). Corollary~\ref{d:numerical} consequently
gives, for nearby data in its sourced-constraint class with weakly future
trapped boundary and fixed axis values,
\begin{equation}
\label{j:local-bound}
 E_{\widetilde g}\geq
 F_*\left(\sqrt{\frac{a(\widetilde g)}{4\pi}}\right).
\end{equation}
The sourced class permits compact variations of the potentials themselves,
provided their polar orders and endpoint constants are retained. For each
such choice of map, the numerical deformation keeps that chosen map fixed.
Thus Equation~\eqref{j:local-bound} applies to all the map variations below;
the conserved jumps \((2Q,0,4J)\) are unchanged. The strict branch
persists by uniform metric equivalence. At the
background, direct differentiation of the model mass gives
\begin{equation}
\label{j:area-coefficient}
 c:=16\pi\frac{\dd}{\dd A}
           F_*\left(\sqrt{\frac{A}{4\pi}}\right)
   =\frac{2b}{R_0^2}=2\kappa>0.
\end{equation}

\subsection{Positive multipliers from the feasible cone}

Let \(\mathcal B\) be the closed unit-ball bundle of \(g\), including the
fibers over \(S\). Transport its vectors isometrically when the metric
varies, and define the residual
\begin{equation}
\label{j:residual}
 C(x,w)=16\pi\bigl(\mu+J_{\rm con}(w)\bigr)
   -2\sum_{i=1}^3\bigl(|D_i|^2+s_i^2+2s_iD_i(w)\bigr),
 \qquad |w|_g\leq1.
\end{equation}
Equation~\eqref{a:constraints} says \(C\equiv0\). Consequently the
infinitesimal contribution from transporting the ball vectors is zero.

The compact variation space contains all smooth axial metric and tensor
variations invariant under azimuthal reflection. Away from the axis,
the potentials and the \(s_i\) vary independently with compact support.
Through an axis we impose their smooth polar orders. In a signed
transverse radius \(y\), the differences of \(\psi,\chi\) from their axis
constants are even and \(O(y^2)\), while
\[
 z+\psi_{\rm ax}\chi-\chi_{\rm ax}\psi
       -z_{\rm ax}=O(y^4)
\]
is even. Here \(\psi_{\rm ax},\chi_{\rm ax},z_{\rm ax}\) are the constants
on the axis component in question. The variables \(s_1,s_2\) are odd
and \(O(y)\); \(s_3\) is even and \(O(y^2)\).
These conditions are preserved by the variations. They follow from
the smooth polar expansions in the axial reduction, including the even
longitudinal and odd transverse components of the axial connection.
For \(i=1,2\), \(D_i\) has an even transverse coefficient and an odd
longitudinal coefficient. The covector \(D_3\) extends smoothly on the
three-dimensional rotational space, with transverse coefficient \(O(y)\).
In particular
\[
 |D_i|^2,\qquad s_i^2,\qquad s_iD_i
\]
are smooth there, so that all residuals and their variations are smooth.

Adjoin two noncompact directions. The first is the energy prototype
\(\dot g=\delta/r\) on the end, with a radial cutoff to zero near \(S\).
Its ADM derivative is \(E'_g=1/2\). The second is the strict conformal
direction from the preparation for Corollary~\ref{d:numerical}. Fix
\(0<\delta<1\), a positive smooth weight \(w_0\) equal to
\(r^{-3-\delta}\) near infinity, and its strictifying function
\(\varphi>0\). Thus
\[
 \partial_\nu\varphi=-1,\qquad
 -\Delta\varphi-|k|\,|\dd\varphi|\geq c_0w_0,\qquad
 \varphi=O_2(r^{-1}),\qquad
 \frac{-\Delta\varphi}{w_0}\longrightarrow1.
\]
The strict direction is the derivative of
\[
 g_d=e^{4d\varphi}g,\qquad k_d=e^{2d\varphi}k,\qquad
 s_{i,d}=e^{-2(1+h_i)d\varphi}s_i,
\]
with fixed potentials. Its derivatives satisfy
\begin{equation}
\label{j:strict-direction}
 \begin{split}
 -\theta_+'&=4\quad\text{on }S,\\
 C'&=-8\Delta\varphi+8k(w,\nabla\varphi)\\
   &\hspace{1.4em}
       +8\varphi\sum_i h_i
                  \bigl(|D_i|^2+s_i^2+2s_iD_i(w)\bigr)>0,
 \qquad \frac{C'}{w_0}\longrightarrow8.
 \end{split}
\end{equation}
Each row in the sum is nonnegative on the unit ball. The energy
prototype has \(C'=O(r^{-4})\): its Euclidean scalar linearization is
\(-2\Delta_\delta(1/r)=0\), and all remaining terms have this decay or
better. Compact directions have zero tail.

\begin{lemma}[Multiplier identity]
\label{j:multipliers}
There are a probability measure \(\omega\) on \(\cT_0\), a nonnegative
finite measure \(\lambda_S\) on \(S\), a nonnegative locally finite
measure \(\Lambda\) on \(\mathcal B\), and \(z_0\geq0\), such that
every variation in the space just described satisfies
\begin{equation}
\label{j:multiplier-identity}
 16\pi E'_g-c\int_{\cT_0} A'_T\,\dd\omega(T)
 =\int_{\mathcal B}C'\,\dd\Lambda
       -\int_S\theta_+'\dd\lambda_S+z_0a_{\rm dir}.
\end{equation}
Here \(a_{\rm dir}\) is the coefficient of the strict direction.
The measures can be chosen invariant under rotations and azimuthal
reflection.
\end{lemma}

\begin{proof}
Let \(\mathcal B^+\) be the one-point compactification of \(\mathcal B\).
For a variation \(v\) consider the continuous test triple
\[
 L(v)=\left(
 \frac{C'[v]}{w_0},\ -\theta_+'[v],\
 \left\{cA'_T[v]-16\pi E'_g[v]\right\}_{T\in\cT_0}
 \right)
\]
in
\[
 \mathcal X=C(\mathcal B^+)\times C(S)\times C(\cT_0).
\]
Continuity at the added point follows from
Equation~\eqref{j:strict-direction} and \(\delta<1\); continuity of the
last block follows from Lemma~\ref{j:area}.

The linear space \(L(V)\) misses the open cone of triples strictly
positive everywhere. Indeed a variation in this cone has a positive
uniform margin in all three blocks. It has a smooth path realizing
its derivative with
\[
 C_d=dC'+O(d^2w_0),\qquad
 \theta_+(d)=d\theta_+'+O(d^2),
\]
uniformly on the indicated domains. For completeness, realize a finite
sum of compact directions, the prototype, and the strict direction by
first making the linear compact and prototype changes, and then applying
the displayed exponential scalings with parameter \(a_{\rm dir}d\).
On a compact set the Taylor remainders are uniform, also through the
axis by the polar orders. On the end the new metric tails are
\(O_2(r^{-1})\); the second-order scalar-curvature remainder is
\(O(d^2r^{-4})\), as are or dominate all other constraint remainders.
Since \(r^{-4}=O(w_0)\), this proves the stated estimate.
The end curvature remains \(O(r^{-3-\delta})\), the constraints remain
integrable, and the ADM energy is differentiable. All other assumptions
of the enlarged class in Corollary~\ref{d:numerical} persist.
For small \(d>0\) the path is therefore feasible. Equation~\eqref{j:envelope}
and the strictly positive objective block make the derivative of
\[
 16\pi\left[
 E_{g_d}-F_*\left(\sqrt{\frac{a(g_d)}{4\pi}}\right)\right]
\]
strictly negative at its zero background value, contradicting
Equation~\eqref{j:local-bound}.

Separation of an open convex cone from a linear subspace gives a
nonzero continuous positive functional on \(\mathcal X\) annihilating
\(L(V)\). This separation requires no closed-range assertion for \(L\).
Riesz representation gives three nonnegative finite measures.
The measure of the objective block has positive total mass: otherwise
evaluation on the strict direction, whose first two blocks are strictly
positive, forces both remaining measures to vanish. Divide by that
positive mass to obtain \(\omega\). If \(\overline\Lambda\) denotes the
resulting first measure, set
\(\dd\Lambda=w_0^{-1}\dd\overline\Lambda\) on \(\mathcal B\).
It is locally finite. Its atom at the added point gives
\(z_0a_{\rm dir}\), since the first test block has value
\(8a_{\rm dir}\) there. This proves
Equation~\eqref{j:multiplier-identity}. Averaging the identity over the
compact symmetry group gives the last assertion.
\end{proof}

\subsection{Concentration on the original boundary}

\begin{lemma}
\label{j:concentration}
The probability measure in Equation~\eqref{j:multiplier-identity}
is concentrated on the cut \(S\).
\end{lemma}

\begin{proof}
Let \(\sigma_0\) be an arbitrary smooth invariant normal speed supported
strictly inside \(\Omega\). Take a short local Einstein--Maxwell Cauchy
development of the \emph{original} data about its compact support.
Local hyperbolic existence and uniqueness propagate the symmetry.
A compact change of its spacelike slice, followed by polarization,
gives exact constraint variations
\begin{equation}
\label{j:slice-test}
 C'=0,\qquad \dot g=2\sigma_0k.
\end{equation}
To justify the second identity, differentiate the orthogonal horizontal
projection defining \(h\). Terms from differentiating the horizontal
lifts pair to zero by orthogonality. Differentiating the fiber length
gives the remaining fiber component, so the polarized variation is
precisely \(2\sigma_0k\).

The potentials in these tests can be chosen with compact variation and
unchanged axis constants. Maxwell's equations give the closed spacetime
orbit forms defining \(\psi,\chi\), which are integrated normally in the
short slab. With these potentials fixed, the orbit form defining \(z\)
is closed as well: its exterior derivative vanishes on every spacelike
meridional plane by the axial momentum equation on each tilted slice.
Such planes test every two-form locally in the three-dimensional orbit
space. Normal integration therefore extends \(z\) with its initial
constants. These constructions preserve the polar orders. Only the
interior neighborhood of the chosen support is used.

All end and boundary variations, and \(a_{\rm dir}\), vanish in these
tests. Equations~\eqref{j:multiplier-identity} and \eqref{j:slice-test} give
\begin{equation}
\label{j:trace-measure}
 \int_{\cT_0}\int_{\partial T}
       \sigma_0\,\tr_{\partial T}k\,\dd A_g\,\dd\omega(T)=0.
\end{equation}
The common-plane property above makes the tangential trace a single
measurable function on the union of free frontiers. The aggregate area
measure is positive. Equation~\eqref{j:trace-measure}, first for invariant
tests and then for all tests by averaging, says that this function
vanishes almost everywhere for the aggregate measure. Integrating its
absolute value and using Fubini shows that, for \(\omega\)-almost every
\(T\), the trace vanishes on almost every point of its free frontier.
Continuity makes it vanish everywhere on each smooth free portion.
Such portions are minimal and hence satisfy the MOTS equation.

The same equation holds almost everywhere across contact with \(S\).
In tangent graphs the difference between the frontier and the obstacle,
and its first derivatives, vanish on the contact set. Since the graphs
are \(C^{1,1}\), differentiability and the density theorem imply that
their second derivatives also agree at almost every contact point.
Their outward orientations agree by obstacle inclusion. The frontier
therefore solves the smooth uniformly elliptic MOTS graph equation
almost everywhere across contact. In a graph chart, write this as
\[
 a^{ij}(x,f,\dd f)\partial_i\partial_jf
                   =b(x,f,\dd f).
\]
The \(C^{1,1}\) graph is a strong \(W^{2,\infty}\) solution; bounded
slope makes the principal matrix uniformly elliptic. Evaluated on
\((x,f,\dd f)\), both the coefficients and the right side are Lipschitz.
Interior Schauder regularity for this linear equation therefore gives
\(f\in C^{2,\alpha}\) for every \(\alpha<1\) on a smaller chart.
Differentiation and elliptic bootstrap then make the graph smooth.
Strong comparison implies local coincidence whenever
it touches \(S\). Connectedness then forces a touching component to be
all of \(S\); its total area is \(A\), so it has no further component.

If the frontier is disjoint from \(S\), it is a smooth interior enclosing
MOTS. For an invariant surface the original and polarized expansions
agree. The filled-frontier convention removes any spurious pocket
component, so this is an enclosing surface forbidden by the exact
outermostness hypothesis of Theorem~\ref{m:main}. Thus almost every
minimizing cut in the multiplier is \(S\).
\end{proof}

\subsection{Smoothness of the lapse and shift}

Take the zeroth and first fiber moments of \(\Lambda\). As distributions
relative to \(\dd V_g\), denote them by \(N\) and \(Y\):
\[
 N\,\dd V_g=\operatorname{pr}_*\Lambda,\qquad
 Y\,\dd V_g=\operatorname{pr}_*(w\Lambda).
\]
Positivity of \(\Lambda\) gives \(N\geq |Y|\) in the measure sense.
Averaging makes \(Y\) meridional. After
Lemma~\ref{j:concentration}, the interior compact tests in
Equation~\eqref{j:multiplier-identity} have only the constraint term.

The variation of the trace-reversed tensor
\(k-(\tr_gk)g\) gives
\begin{equation}
\label{j:shift}
 \operatorname{sym}\nabla Y=-Nk.
\end{equation}
For example, if \(P=\dot k-(\tr_g\dot k)g\), then in dimension three
\(\dot k=P-\tfrac12(\tr_gP)g\). The tensor-dependent variation after
integration by parts is
\(-2\langle Nk+\operatorname{sym}\nabla Y,P\rangle\), which proves
Equation~\eqref{j:shift}.
Compact conformal tests with the same source weights as in
Equation~\eqref{j:strict-direction} give
\begin{equation}
\label{j:lapse}
 \Delta N=-\divg k(Y,\cdot)
   +\sum_i h_i\left[N(|D_i|^2+s_i^2)+2s_iD_i(Y)\right].
\end{equation}
These distributional equations hold through the axis: their coefficients
are smooth, and averaging arbitrary test components gives precisely the
invariant tests already allowed.

\begin{lemma}[Regularity and a first-jet system]
\label{j:regularity}
The interior moment distributions are smooth. Their first jets satisfy
a homogeneous linear first-order system with coefficients smooth through
the axis and one-sided smooth up to \(S\). They extend smoothly to \(S\),
and the moment measures have no part supported on \(S\).
\end{lemma}

\begin{proof}
Take the divergence of Equation~\eqref{j:shift} together with
Equation~\eqref{j:lapse}. The scalar principal symbol is \(|\xi|^2\);
the vector symbol is
\[
 \frac12\bigl(|\xi|^2I+\xi\otimes\xi\bigr).
\]
All couplings are lower order. This is an elliptic system with smooth
coefficients on the rotational three-space. Distributional elliptic
regularity therefore proves interior smoothness, including at the axis.

We record the additional metric equation to justify continuation and the
end estimates. For a pure metric variation the second-order term is
\(\Hess N-(\Delta N)g\). The gravitational lower-order terms are smooth
linear expressions bounded in norm by
\begin{equation}
\label{j:gravitational-coefficients}
 C\left[(|\operatorname{Rm}|+|k|^2)|N|
             +|\nabla k|\,|Y|+|k|\,|\nabla Y|\right].
\end{equation}
This follows directly from
\[
 R'=\divg\divg\dot g-\Delta\tr_g\dot g-\langle\Ric,\dot g\rangle,
\]
the algebraic variation of \((\tr k)^2-|k|^2\), and the
\(\nabla k*\dot g+k*\nabla\dot g\) terms in the momentum variation.
Terms due to varying the volume disappear because the background
constraint residual is zero.

The matter terms, while initially written off the axis, have smooth
tensor coefficients. Up to constant factors and smooth trace terms,
they contract against
\begin{equation}
\label{j:matter-coefficients}
 N\left(D_i^\sharp\otimes D_i^\sharp+
                      h_i|D_i|^2e\otimes e\right),
       \qquad s_iD_i(Y)e\otimes e.
\end{equation}
Free compact variations of \(s_i\) give
\begin{equation}
\label{j:map-velocity}
 D_i(Y)=-Ns_i.
\end{equation}
Thus the second tensors become \(-Ns_i^2e\otimes e\), which are
smooth by the polar orders. For \(i=1,2\), the radial and azimuthal
coefficients of the first tensor have the same leading value on the axis
and smooth even parity; their difference is \(O(y^2)\). The
radial-longitudinal coefficient is smooth odd. In Cartesian polar disks
these are exactly the conditions for a smooth invariant tensor.
For \(i=3\), the higher vanishing orders directly give the same conclusion.
This proves smoothness of the full metric equation through the axis; the
matter terms are bounded by
\(C|N|\sum_i(|D_i|^2+s_i^2)\).

Taking the trace of the metric equation and substituting back determines
all components of \(\Hess N\) from the first jets of \(N,Y\).
Differentiating Equation~\eqref{j:shift} and commuting covariant derivatives
likewise expresses \(\nabla^2Y\) in those first jets, with coefficients
formed from \(k,\nabla k,\operatorname{Rm}\). Hence
\[
 (N,Y,\nabla N,\nabla Y)
\]
satisfies the asserted homogeneous first-order system. In a collar of
\(S\), integrate its normal ordinary differential equation from an
interior collar slice. Smooth dependence on the tangential parameters
gives a smooth one-sided extension to the boundary.

There might initially remain a separate boundary-supported part of the
moment measures. Take \(\dot k=0\) and a compact metric variation whose
value and first jet vanish on \(S\), with arbitrary second normal
derivative of its tangential trace there. Neither the area nor the
expansion derivative sees this second jet. After integration of the
smooth interior adjoint, its boundary terms also vanish. In the residual
at \(S\), only the scalar-curvature variation tests the prescribed
second jet. Equation~\eqref{j:multiplier-identity} therefore forces the
boundary part of \(N\) to vanish. The measure inequality \(|Y|\leq N\)
then removes the boundary part of \(Y\).
\end{proof}

\subsection{End normalization}
\label{j:end-section}

\begin{lemma}[Asymptotic constants and the spatial translation]
\label{j:end}
For the adjoint fields above, there are constants \(N_\infty,Y_\infty\)
such that, for every \(q'\in(1/2,1)\),
\[
 (N,Y)=(N_\infty,Y_\infty)+O_2(r^{-q'}).
\]
The multiplier identity gives \(N_\infty=1\), and positive ADM energy
gives \(Y_\infty=0\). In fact \(Y=O_2(r^{-2})\).

More generally, the elimination of \(Y_\infty\) applies to any pair
with these constant asymptotics that satisfies
\(\operatorname{sym}\nabla Y=-Nk\), provided
\(k=O_1(r^{-3})\), \(N\) is bounded, the scalar linearization
\(L(g-\delta)=O(r^{-4})\), and the ADM energy is positive.
\end{lemma}

\begin{proof}
Set \(F=(N,Y)\). The first-jet equations and the original end bounds give
\begin{equation}
\label{j:end-ode}
 |\partial^2F|
 \leq C r^{-2-q'}|F|+C r^{-1-q'}|\partial F|.
\end{equation}
Here the matter coefficients have order at least \(r^{-4}\), so any
\(q'\in(1/2,1)\) is allowed. Radial integration with Gronwall's inequality,
applied to \(|F|/r+|\partial F|\), first gives
\(F=O(r)\) and \(\partial F=O(1)\).
The Hessian bound then gives a limiting gradient along each ray.
Integrating the Hessian along paths of length \(O(r)\) on coordinate
spheres shows that these limits agree. A nonzero linear term in \(N\)
would contradict \(N\geq0\) in opposite asymptotic directions. Once that
term vanishes, \(|Y|\leq N\) eliminates every linear term of \(Y\) as
well. Consequently
\[
 F=O(r^{1-q'}),\qquad \partial F=O(r^{-q'}).
\]
Reinsertion into Equation~\eqref{j:end-ode} gives
\(\partial^2F=O(r^{-1-2q'})\). The gradient limits are zero, so
\(\partial F=O(r^{-2q'})\). Since \(2q'>1\), radial integration gives
finite value limits, and sphere paths show that there is a single such
limit. One final insertion gives
\[
 \partial^2F=O(r^{-2-q'}),\qquad
 \partial F=O(r^{-1-q'}),\qquad
 F-F_\infty=O(r^{-q'}).
\]

Use the energy prototype in Equation~\eqref{j:multiplier-identity}.
It vanishes near \(S\), and \(a_{\rm dir}=0\). Integrate the smooth
adjoint over a large truncation. Scalar linearization gives the limiting
constraint flux \(16\pi N_\infty E'_g\). Terms containing \(\nabla N\),
the nonconstant part of \(N\), and momentum metric-variation boundary
terms tend to zero by the preceding decay. The residual tail is
integrable. Since \(E'_g=1/2\neq0\), the identity gives \(N_\infty=1\).

To eliminate the constant spatial translation, write
\(g=\delta+H_0\) and put \(a=Y_\infty\) for this paragraph only.
Equation~\eqref{a:constraints} and scalar linearization give
\[
 LH_0:=\partial_i\partial_jH_{0,ij}
                   -\Delta_\delta\tr_\delta H_0=O(r^{-4}).
\]
For an affine function \(V\) define the Euclidean flux covector
\[
 U(V;h')_i=
 V(\partial_jh'_{ij}-\partial_ih'_{jj})
       -h'_{ij}\partial_jV+(\tr_\delta h')\partial_iV.
\]
It satisfies \(\partial_iU(V;h')_i=VLh'\).
The shift equation gives \(\cL_Yg=O_1(r^{-3})\). Subtracting
\(\partial_aH_0\) and the pure Euclidean gauge \(\cL_Y\delta\)
leaves errors \(O_1(r^{-2-q'})\). The affine flux of these errors tends
to zero. The affine flux of \(\cL_Y\delta\) is identically zero after a
smooth extension to the interior, because its scalar linearization
vanishes. Thus the flux of \(U(V;\partial_aH_0)\) tends to zero.

On the other hand, differentiation of the definition of \(U\) gives
\[
 U(V;\partial_aH_0)=\partial_aU(V;H_0)-U(\partial_aV;H_0).
\]
Euclidean integration by parts, with arbitrary smooth interior
extensions, consequently yields the exact sphere identity
\begin{equation}
\label{j:boost-flux}
 \begin{split}
 \int_{S_r}U(V;\partial_aH_0)\cdot n_\delta\,\dd A_\delta
 &=\int_{S_r}(a\cdot n_\delta)V\,LH_0\,\dd A_\delta\\
 &\quad-(\partial_aV)
         \int_{S_r}U(1;H_0)\cdot n_\delta\,\dd A_\delta.
 \end{split}
\end{equation}
The first term tends to zero; the second tends to
\(-16\pi m\,\partial_aV\). Since \(m>0\), testing all affine \(V\)
gives \(a=0\). This calculation uses just the hypotheses stated in the
last part of the lemma.

Finally the prolonged shift equation, now with zero limits, gives
\[
 |\partial^2Y|
 \leq C\bigl(r^{-2}|\partial Y|+r^{-3}|Y|+r^{-4}\bigr).
\]
Starting with \(Y=O_2(r^{-q'})\), integration from infinity first
improves this to \(Y=O_2(r^{-1-q'})\), and then to
\(Y=O_2(r^{-2})\).
\end{proof}

We also obtain strict interior positivity. If \(N\) vanished at an interior
point, its nonnegativity would give \(N=O(\dist^2)\) there. Domination
would give the same bound for \(Y\). Thus the entire first jet of \(N,Y\)
would vanish at that point. Uniqueness for the homogeneous first-jet
ordinary differential equation along paths would force both fields to
vanish throughout the connected interior, contradicting \(N_\infty=1\).

\subsection{Boundary data and surface gravity}

Compact tensor variations up to \(S\), integrated in
Equation~\eqref{j:multiplier-identity}, give
\begin{equation}
\label{j:boundary-measure}
 Y_T=0,\qquad \dd\lambda_S=2Y_\nu\,\dd A_g.
\end{equation}
Indeed the boundary term of the momentum pairing, whose integration
normal is \(-\nu\), is
\(-2Y^i(\dot k_{ij}-(\tr\dot k)g_{ij})\nu^j\).
Its mixed tangential-normal components force \(Y_T=0\), while its
tangential trace component gives the measure identity. Axial and
reflection symmetry supply all components needed for the meridional
field \(Y\).

Now use the weighted conformal variation with an arbitrary compact
smooth scalar \(\varphi_0\), allowing independent values and normal
derivatives on \(S\). The area derivative is
\(4\int_S\varphi_0\,\dd A_g\), and
\(\theta_+'=4\partial_\nu\varphi_0\).
Integrating Equation~\eqref{j:lapse}, with integration normal \(-\nu\),
gives
\begin{equation}
\label{j:boundary-conformal}
 \begin{split}
 -4c\int_S\varphi_0\,\dd A_g
 &=8\int_S\left[
       N\partial_\nu\varphi_0
       -\bigl(\partial_\nu N+k(Y,\nu)\bigr)\varphi_0
             \right]\dd A_g\\
 &\quad-4\int_S\partial_\nu\varphi_0\,\dd\lambda_S.
 \end{split}
\end{equation}
Independent comparison of the two boundary jets, together with
Equation~\eqref{j:boundary-measure}, gives
\(N=Y_\nu\) and
\(\partial_\nu N+k(Y,\nu)=c/2\).
Equation~\eqref{j:area-coefficient} proves
Equation~\eqref{j:boundary} with its stated positive surface gravity.

\begin{lemma}[Boundary lapse dichotomy]
\label{j:dichotomy}
Either \(N|_S>0\) everywhere or \(N|_S\equiv0\).
\end{lemma}

\begin{proof}
Choose a compact smooth invariant meridional vector field \(W\) with
\(W|_S=\xi_0\nu\). Extend the data smoothly across \(S\) and take their
Lie derivatives along \(W\). These are allowed compact first variations
in Equation~\eqref{j:multiplier-identity}; they need not form a feasible
path. Covariance gives \(C'=\cL_WC=0\) on the exterior. This also
holds at \(S\), since the exact residual vanishes on the interior and
all its one-sided jets vanish at the boundary. No constraint condition
on the extension is required. Naturality identifies the fixed-surface
expansion of the pulled-back data with the expansion of the surface
moved by \(W\) in the original data. Consequently the first area
variation is \(\int_S H\xi_0\,\dd A_g\), and the expansion variation is
\(L_S\xi_0\), where \(L_S\) is the MOTS normal stability operator with
principal part \(-\Delta_S\). Equation~\eqref{j:multiplier-identity} and
\(\dd\lambda_S=2N\,\dd A_g\) therefore give
\begin{equation}
\label{j:stability-boundary}
 2L_S^*N=cH.
\end{equation}
Averaging allows arbitrary scalar tests. Outer area minimization implies
\(H\geq0\) by outward first variation; the polarized and original mean
curvatures of \(S\) agree. Thus \(L_S^*N\geq0\) with \(N\geq0\).
To check the sign in the strong minimum principle, write locally
\(L_S^*=-a^{ij}\partial_i\partial_j+b^i\partial_i+c_0(x)\).
Choose a constant \(C_0\geq\max(0,\sup_Sc_0)\). Replacing \(c_0\)
by \(C_0\) preserves the inequality because \(N\geq0\). The function
\(v=-N\leq0\) then satisfies
\[
 (a^{ij}\partial_i\partial_j-b^i\partial_i-C_0)v\geq0.
\]
Its zeroth-order coefficient is nonpositive, and a zero of \(N\) is a
nonnegative interior maximum of \(v\). The strong maximum principle
forces local constancy. Connectedness of \(S\) proves the dichotomy.
\end{proof}

The multiplier at infinity can now also be identified. Apply the smooth
conformal adjoint identity to the strictifying function \(\varphi\),
integrating over larger truncations. Its source terms are integrable,
and the end flux is \(16\pi E'_g\) because \(N_\infty=1\); all
remaining end terms vanish by Equation~\eqref{j:asymptotic}. The boundary
terms are exactly those of Equation~\eqref{j:boundary-conformal}.
Thus
\[
 \int_{\mathcal B}C'\,\dd\Lambda
      -\int_S\theta_+'\dd\lambda_S
       =16\pi E'_g-cA'_S
\]
for the strict direction. Comparing with
Equation~\eqref{j:multiplier-identity}, where \(a_{\rm dir}=1\), gives
\(z_0=0\). Its value was not needed earlier: every compact test and
the energy prototype have \(a_{\rm dir}=0\).

\subsection{The canonical quotient action}

We finish the proof of Proposition~\ref{j:adjoint} by interpreting the
complete interior adjoint. Work off the axis, where \(X>0\), and introduce
the quotient canonical variables
\begin{equation}
\label{j:canonical-variables}
 \widehat q=X^2h,\qquad m_v=k(e,e),\qquad
 p=X(k|_h+m_vh).
\end{equation}
Let \(\pi\) be the target tangent vector with orthonormal components
\begin{equation}
\label{j:canonical-momentum}
 \dd u(\pi)=X^{-1}m_v,\qquad
 \cD_X(\pi)=X^{-1}s.
\end{equation}
Together with \(\Phi\), the variables \(\widehat q,p,\pi\) are free local
variation coordinates. The quotient lapse is \(\widehat N=XN\).
With all spatial contractions now taken using \(\widehat q\), define
\begin{equation}
\label{j:canonical-constraints}
 \begin{split}
 \widehat C_H
 &=R_{\widehat q}+(\tr_{\widehat q}p)^2-|p|_{\widehat q}^2
                  -2\bigl(|\dd\Phi|_{\widehat q,G}^2+|\pi|_G^2\bigr),\\
 \widehat C_P
 &=\divg_{\widehat q}
           \bigl(p-(\tr_{\widehat q}p)\widehat q\bigr)
                      -2G(\pi,\dd\Phi).
 \end{split}
\end{equation}
The exact relation to the polarized residual is
\begin{equation}
\label{j:constraint-conversion}
 C(x,w)=X^2\widehat C_H+2X\widehat C_P(w).
\end{equation}
Here only the horizontal component of \(w\) enters the second term;
the axial momentum residual vanishes.

We verify both parts of this conversion. The scalar identities are
\[
 X^2R_{\widehat q}=R_g+2|\dd\log X|_h^2,
 \qquad
 X^2\bigl((\tr p)^2-|p|^2\bigr)-2m_v^2
                      =(\tr_gk)^2-|k|_g^2.
\]
They follow respectively from two-dimensional conformal curvature and
the trace algebra of Equation~\eqref{j:canonical-variables}. The spatial
and normal \(u\)-energies cancel the two extra terms. For momentum, put
\[
 P_H=k|_h-(\tr_h k|_h+m_v)h.
\]
Its three-dimensional horizontal divergence is
\[
 \divg_hP_H+k|_h(\nabla\log X,\cdot)-m_v\dd\log X.
\]
The rescaled quotient divergence is
\[
 X\divg_{\widehat q}(XP_H)
   =\divg_hP_H+k|_h(\nabla\log X,\cdot)+m_v\dd\log X.
\]
Subtracting the \(u\)-momentum term \(2m_v\dd\log X\), and then the
remaining map momenta, gives the required horizontal residual. This
proves Equation~\eqref{j:constraint-conversion}.

After circle integration the volume identity is
\(\dd V_g=2\pi X^{-1}\dd A_{\widehat q}\). Since the constraints vanish
at the background, variations of this conversion factor or of the
multiplier coordinates make no contribution. The compact adjoint
identity is therefore stationarity of
\begin{equation}
\label{j:canonical-functional}
 I=\int\left[
      \widehat N\widehat C_H+2Y\cdot\widehat C_P
          \right]\dd A_{\widehat q}
\end{equation}
under all compact variations of
\(\widehat q,p,\Phi,\pi\). Variations of \(\widehat N,Y\) themselves
also vanish by the constraints.

Integrating the momentum divergence by parts and varying \(p,\pi\)
gives, respectively,
\begin{equation}
\label{j:stationary-momenta}
 p=-\frac{\cL_Y\widehat q}{2\widehat N},
 \qquad
 \pi=-\frac{\dd_Y\Phi}{\widehat N}.
\end{equation}
For instance the \(p\)-equation before tracing is
\[
 \widehat N\bigl((\tr p)\widehat q-p\bigr)
 -\operatorname{sym}\widehat\nabla Y
            +(\divg_{\widehat q}Y)\widehat q=0.
\]
Its two-dimensional trace eliminates
\(\widehat N\tr p+\divg_{\widehat q}Y\), giving the first equation.
The \(\pi\)-variation is
\(-4\widehat N\pi-4\dd_Y\Phi=0\), giving the second.

Since \(N>0\) in the interior, these are precisely the induced tensor
and normal map velocity for the Lorentzian lapse-shift metric
\[
 \gamma=-\widehat N^2\dd t_0^2+
  \widehat q_{ij}(\dd x^i+Y^i\dd t_0)(\dd x^j+Y^j\dd t_0)
\]
with stationary coefficients and the positive second-fundamental-form
convention. Substitution into Equation~\eqref{j:canonical-functional}
gives
\begin{equation}
\label{j:stationary-action}
 \int\widehat N\left[
 R_{\widehat q}+|p|^2-(\tr p)^2
           -2|\dd\Phi|_{\widehat q,G}^2+2|\pi|_G^2
                        \right]\dd A_{\widehat q}.
\end{equation}
The signs can also be read directly from the integrated shift terms:
they add \(2\widehat N(|p|^2-(\tr p)^2)\) and
\(4\widehat N|\pi|^2\) to the original constraint expression.

The Gauss decomposition identifies
Equation~\eqref{j:stationary-action}, up to divergences, with
\[
 \int\bigl(R_\gamma-2|\dd\Phi|_\gamma^2\bigr)\dd V_\gamma
\]
per unit \(t_0\)-time. Thus its metric and map variations give
\[
 \Ric_\gamma=2\Phi^*G,\qquad
 \operatorname{tr}_\gamma\nabla\dd\Phi=0,
\]
which is Equation~\eqref{a:wave}. Stationary tests give all equations:
one may take a periodic time interval for the compact variational
calculation and average any test in time, since every coefficient and
Euler--Lagrange expression is stationary.

Finally the \(H_3\)-normal is \(n=Xn_\gamma\). Equations
\eqref{j:canonical-momentum} and \eqref{j:stationary-momenta} therefore
give \(\dd u(n)=m_v\) and \(\cD_X(n)=s\).
Under \(\gamma=X^2H_3\) the slice tensors are related by
\[
 p=X\bigl(k_{H_3}+n(\log X)h\bigr),
\]
so Equation~\eqref{j:canonical-variables} gives \(k_{H_3}=k|_h\).
Together with the preceding lemmas, this proves
Proposition~\ref{j:adjoint}.
